% ============================================================
%  THE ORTYX PROTOCOL
%  Part III — Decomposition
%  Chapter 15: The Audit Summary and Preparation
%              for the Interlude
% ============================================================

\chapter{The Audit Summary and Preparation for the Interlude}
\label{ch:audit-summary}

\chapepigraph{Decomposition is not destruction. It is
the separation of components that were fused together
without argument. Some of what is separated will prove
to be stronger when standing alone than it appeared
when supporting claims it could not carry.}{Flyxion}

\noindent
Part~III has conducted a systematic formal audit of
the Phoenix corpus using the four audit operators
of the \Ortyx{} Protocol. The results are now assembled
into a summary that prepares the Interlude's explicit
classification of survivable and unstable components.

\section{The Phoenix Audit Score}
\label{sec:phoenix-score}

Applying the epistemic stability functional with weights
reflecting the Phoenix corpus's status as an early-stage
research programme ($\alpha = 0.3, \beta = 0.4,
\gamma = 0.2, \delta = 0.1$, reflecting the
importance of falsifiability at this stage):

\[
  \EpStab(\SysTriple_{\mathrm{Phoenix}})
  = 0.3\,\AuditS + 0.4\,\AuditF
  - 0.2\,\AuditP - 0.1\,\AuditR.
\]

Estimated operator values, disaggregated by claim level:

\begin{center}
\begin{tabular}{lcccc}
\toprule
Claim Level & $\AuditS$ & $\AuditF$ & $\AuditP$ & $\AuditR$ \\
\midrule
\LevelI{} (Engineering) & 0.85 & 0.80 & 0.25 & 0.20 \\
\LevelII{} (Metaphor)   & 0.65 & 0.55 & 0.45 & 0.35 \\
\LevelIII{} (Phenomenol.) & 0.50 & 0.40 & 0.60 & 0.45 \\
\LevelIV{} (Ontological) & 0.25 & 0.15 & 0.80 & 0.70 \\
\bottomrule
\end{tabular}
\end{center}

Applying the functional at each level:
\[
  \EpStab^{(I)} \approx 0.71, \quad
  \EpStab^{(II)} \approx 0.41, \quad
  \EpStab^{(III)} \approx 0.22, \quad
  \EpStab^{(IV)} \approx 0.01.
\]

The gradient is striking: from a respectable score
at \LevelI{} to near-zero at \LevelIV{}. This confirms
the structural diagnosis of Part~II: the Phoenix corpus
is well-grounded at the engineering level and loses
epistemic stability progressively as claims ascend
the level hierarchy.

\begin{epistemicstatus}{Operationally Grounded}
The audit score values presented above are estimates
rather than precise calculations, for the reason
identified in Chapter~11: the operator values require
human judgment about what constitutes falsification,
constraint-faithfulness, and closure, and these
judgments are not uniquely determined by algorithm.
The values represent the author's considered assessment
after the analysis of Chapters~11--14; they should
be read as order-of-magnitude estimates with significant
uncertainty, not as exact measurements. The gradient
--- the substantial drop from \LevelI{} to \LevelIV{} ---
is more reliable than the specific values.
\end{epistemicstatus}

\section{The Four Failure Geometries: Summary}
\label{sec:fg-summary}

\begin{center}
\begin{tabular}{lll}
\toprule
Geometry & Primary Locus & Severity \\
\midrule
$\FailGeo{1}$ Decorative Formalism & $H_i$ in Marine; $e(t)$ in \MEMeight{} & Low--Moderate \\
$\FailGeo{2}$ Definitional Closure & Authenticity; salience--meaning & Moderate \\
$\FailGeo{3}$ Semantic Universality & Coherence cascade across all levels & High \\
$\FailGeo{4}$ Operational Severance & \texttt{.m8a} consciousness metadata & High \\
\bottomrule
\end{tabular}
\end{center}

\section{What Has Been Established}
\label{sec:what-established}

Part~III has established the following:

The Phoenix corpus contains genuine engineering results
that survive the full audit at \LevelI{}: the jitter
metrics, the conformity relation, the associative
memory efficiency of wave superposition, and the
context-sensitivity vulnerability of LLMs.

The Phoenix corpus contains productive computational
metaphors at \LevelII{} that survive the audit as
research heuristics: the salience--accessibility
entropy connection, the conformity--population coding
analogy, the memory-interference parallel, and the
behavioral injection framing.

The Phoenix corpus contains phenomenologically grounded
observations at \LevelIII{} that survive the audit
independently of the mechanisms: the temporal character
of musical experience, the reconstruction character
of memory, and the phenomenology of compression
impoverishment.

The Phoenix corpus's ontological claims at \LevelIV{}
do not survive the audit at their full strength.
The consciousness claims, the authenticity ontology,
and the strong AI architecture thesis are in $\ThetaU$
after decomposition.

The RSVP framework passes the audit at a level
sufficient to serve as a reinterpretive substrate,
with explicit acknowledgments of its own vulnerabilities.

\section{What Has Not Been Established}
\label{sec:what-not-established}

Part~III has not established that the Phoenix corpus
is false. The ontological claims at \LevelIV{} may
be correct; the audit does not determine this. What
it determines is that they are not established by
the framework's current apparatus, and that their
current formulation exhibits failure geometries
that prevent them from being evaluated as scientific
claims.

Part~III has not established that the failure geometries
are unique to the Phoenix corpus. On the contrary,
Chapter~11 noted that these geometries are characteristic
of early-stage speculative frameworks generally, and
Chapter~14 found that RSVP itself is not free of
the risks associated with $\FailGeo{3}$ and
projection dependence.

Part~III has not determined the weights
$\alpha, \beta, \gamma, \delta$ in any principled
way beyond the context-sensitive judgment described
in Chapter~11. Different weight choices would yield
different audit scores while preserving the gradient.

\section{The Preparation}
\label{sec:preparation}

The Interlude that follows Part~III will perform the
explicit classification of components into the five
categories required for the reconstruction work of Part~IV:

\textit{Discarded}: Components that are not recoverable
within any reasonable reinterpretation, because they
exhibit $\FailGeo{2}$ or $\FailGeo{3}$ at severity
levels that prevent them from making genuine empirical
contact.

\textit{Unstable but recoverable}: Components that
are currently in $\ThetaU$ due to missing specification
or operationalization work, but that could migrate
to $\ThetaS$ if the specification gaps were closed.
These are the most valuable targets for future
framework development.

\textit{Metaphorically useful}: Components that
function productively as \LevelII{} analogies and
research heuristics, even though they cannot bear
the weight of \LevelIV{} claims. These survive as
productive conceptual tools without claiming to
be established mechanisms.

\textit{Operationally grounded}: Components that
are fully specified, testable, and in $\ThetaS$ at
\LevelI{}. These survive as the engineering core.

\textit{Reconstructible}: Components that are
currently in $\ThetaU$ but that can be formally
translated into a more stable semantic framework
--- specifically, RSVP --- without loss of their
essential content. These are the targets for the
reinterpretation functor of Part~IV.
\[
  \RSVPfunctor : \ThetaS \to \ThetaStar.
\]

This functor does not change what the survivable
components say. It changes the vocabulary in which
they are expressed, replacing vocabulary that carries
projection risks with vocabulary that makes its
projection dependencies explicit. The reinterpretation
is not a claim that RSVP is true and Phoenix is false;
it is a claim that the Phoenix corpus's survivable
components are better expressed in RSVP-compatible
terms because those terms make the framework's
assumptions more auditable.

\medskip

The Interlude follows immediately. It is a different
kind of text than the chapters that surround it:
briefer, more direct, and deliberately transitional.
Its function is to stand at the threshold between
decomposition and reconstruction --- to name what
has been lost and what remains, before the work
of reconstruction begins.
